\section{Request Validator}
\subsection{Overview}

The Request Validator is a purely combinational block instantiated once
each on the read and write paths. It pops entries from the AR or AW
Metadata FIFO, performs a single runtime check (address range
validation), and steers the transaction to one of two downstream paths:

\begin{itemize}
  \item \textbf{Valid path:} Transaction forwarded to ROB allocation
        with \texttt{REQ\_VALID} asserted.
  \item \textbf{Error path:} Transaction forwarded to the Error Handler
        with \texttt{ERR\_VALID} asserted and \texttt{AXID} preserved.
\end{itemize}

The Request Validator contains a Configuration Register that stores the valid physical address range used for runtime address validation. 
Apart from this internal configuration state, the Request Validator contains no per-transaction storage or counters.
All protocol-level checks (burst type, WRAP legality, alignment, burst
length) are enforced as design-time assertions in the verification
environment only. Compliant AXI4 masters are assumed at runtime.

% -----------------------------------------------------------------------
\subsection{Design Parameters}

\begin{table}[h!]
\centering
\begin{tabular}{lll}
\toprule
\textbf{Parameter} & \textbf{Value} & \textbf{Description} \\
\midrule
\texttt{ADDR\_W}  & 32 & Address width in bits \\
\texttt{AXID\_W}  & 6  & AXI ID width \\
\texttt{LEN\_W}   & 8  & Burst length field width (ARLEN/AWLEN) \\
\bottomrule
\end{tabular}
\caption{Request Validator parameters}
\end{table}

% -----------------------------------------------------------------------
\subsection{Internal Blocks}

Configuration Register
Stores the valid physical address range used during request validation.
Provides the configured address range to the Address Range Check block.
Represents the only persistent state within the Request Validator.

\subsection{Address Range Check}

\begin{itemize}
  \item \textbf{Type:} Purely combinational.
  \item \textbf{Input:} \texttt{ADDR[31:0]} from the AR/AW Metadata
        FIFO.
  \item \textbf{Function:} Checks whether \texttt{ADDR} falls within
        the valid physical address range stored in the Configuration Register.
  \item \textbf{Output:} \texttt{addr\_ok} (1-bit) — asserted when the
        address is valid, deasserted on an address error.
  \item \textbf{Scope:} Address range check only. No burst type, no
        alignment, no length checks at runtime.
\end{itemize}

\subsection{Output Mux}

\begin{itemize}
  \item \textbf{Type:} Purely combinational.
  \item \textbf{Function:} Routes the transaction to the correct
        downstream block based on \texttt{addr\_ok}.
  \item \textbf{Valid path} (\texttt{addr\_ok == 1}): asserts
        \texttt{REQ\_VALID} to ROB allocation; forwards
        \texttt{AXID[5:0]}.
  \item \textbf{Error path} (\texttt{addr\_ok == 0}): asserts
        \texttt{ERR\_VALID} to the Error Handler; forwards
        \texttt{AXID[5:0]}.
  \item Only one output path is active per transaction. Both paths are
        never asserted simultaneously.
\end{itemize}

% -----------------------------------------------------------------------
\subsection{Interface Signals}

\subsection{Inputs}

\begin{longtable}{llll}
\toprule
\textbf{Signal} & \textbf{Width} & \textbf{Source} & \textbf{Description} \\
\midrule
\texttt{ADDR[31:0]}        & 32 & AR/AW Metadata FIFO & Transaction address \\
\texttt{AXID[5:0]}         & 6  & AR/AW Metadata FIFO & AXI ID \\
\texttt{ARLEN/AWLEN[7:0]}  & 8  & AR/AW Metadata FIFO & Burst length \\
\texttt{STALL}             & 1  & ROB Allocation      & Per-AXID stall; holds output until deasserted \\
\bottomrule
\caption{Request Validator input signals}
\end{longtable}

\subsection{Outputs}

\begin{longtable}{llll}
\toprule
\textbf{Signal} & \textbf{Width} & \textbf{Destination} & \textbf{Description} \\
\midrule
\texttt{REQ\_VALID}       & 1 & ROB allocation & Valid transaction request \\
\texttt{AXID[5:0]}        & 6 & ROB allocation & Forwarded AXI ID \\
\texttt{ERR\_VALID}       & 1 & Error Handler & Address error notification \\
\texttt{AXID[5:0]}        & 6 & Error Handler & AXI ID preserved for error-response handling \\
\bottomrule
\caption{Request Validator output signals}
\end{longtable}

% -----------------------------------------------------------------------
\subsection{Functional Behaviour}

\begin{enumerate}
  \item Request Validator pops the head entry from the AR/AW Metadata
        FIFO: \texttt{ADDR[31:0]}, \texttt{AXID[5:0]},
        \texttt{ARLEN/AWLEN[7:0]}.
  \item Address Range Check compares the incoming address against the physical address range stored in the Configuration Register and evaluates \texttt{addr\_ok} combinationally.
  \item If \texttt{addr\_ok == 1} and \texttt{STALL} is deasserted:
        Output Mux asserts \texttt{REQ\_VALID} to ROB allocation
        with \texttt{AXID}.
  \item If \texttt{addr\_ok == 0}: Output Mux asserts \texttt{ERR\_VALID}
        to the Error Handler with \texttt{AXID}. The Error Handler sets
        \texttt{pkt\_type = Err\_write / Err\_read} and injects the
        transaction into the normal pipeline via ROB allocation.
  \item If \texttt{STALL} is asserted (ROB Allocation unavailable for this
        AXID): the Request Validator holds its outputs stable and does
        not pop the next FIFO entry until \texttt{STALL} deasserts.
\end{enumerate}

% -----------------------------------------------------------------------
\subsection{Timing Notes}

\begin{itemize}
  \item Fully combinational. Output (\texttt{REQ\_VALID} or
        \texttt{ERR\_VALID}) is valid in the same cycle as the FIFO
        head entry is presented.
  \item \texttt{STALL} from ROB allocation is also combinational —
        the combined path (FIFO output → Address Range Check → Output
        Mux → ROB Allocation Stall Logic → STALL back) is a
        combinational loop that must be timed carefully. A register
        stage may be inserted between the Request Validator output and
        the ROB Allocation if timing closure requires it, at the
        cost of one cycle of allocation latency.
  \item The Address Range Check remains purely combinational. 
        The Configuration Register provides the address range used during validation and is treated as stable during normal request processing.
\end{itemize}

% -----------------------------------------------------------------------
\subsection{Edge Cases}

\begin{itemize}
  \item \textbf{Error path stall:} An address error transaction still
        consumes a ROB allocation. If the AXID is at maximum
        outstanding (\texttt{STALL} asserted), the error path is also
        held. \texttt{ERR\_VALID} is not asserted until \texttt{STALL}
        deasserts.
  \item \textbf{Protocol violations:} Illegal burst type, unaligned
        WRAP, and maximum burst length violations are not checked at
        runtime. They are enforced as design-time assertions in the
        verification environment. A non-compliant master will cause
        undefined behaviour.
  \item \textbf{Simultaneous valid and error:} Not possible —
        \texttt{addr\_ok} is a single-bit result; exactly one output
        path is active per transaction.
\end{itemize}